\chapter{Density functional theory}
\label{chap:dft}
%=============================================================================

Every method so far has been a method for the wave function.  Full
configuration interaction stores its coefficients, Hartree-Fock optimises the
orbitals it is built from, and the random-phase approximation describes the
fluctuations around that.  All of them run into the same wall: the object
being computed lives in a space whose dimension grows exponentially with the
number of particles, Section~\ref{sec:exponentialwall}.

Density functional theory takes a different route.  It replaces the wave
function of $N$ electrons -- a function of $3N$ coordinates -- by the
electronic density, a function of three.  This is not an approximation but,
remarkably, an exact reformulation: Hohenberg and Kohn proved that the
ground-state density determines the external potential, hence the Hamiltonian,
hence everything.  The catch, which we shall be careful to state precisely, is
that the theorem asserts the existence of a universal functional without
telling us its form.

The chapter has three parts.  We first assemble the tools -- the
Born-Oppenheimer Hamiltonian and the reduced density matrices of
Sections~\ref{sec:naturalorbitals} and~\ref{sec:energyfunctional}, now in
coordinate space.  We then prove the Hohenberg-Kohn theorems and examine what
they do and do not deliver.  Finally we build the Kohn-Sham scheme, which
recovers most of what the mean-field methods of
Chapter~\ref{chap:hartrefock} do well while replacing the expensive part by a
local functional, and we construct that functional explicitly from the
electron gas of Section~\ref{sec:electrongas}.

Density functional theory has been an extraordinary practical success.  Most
crystalline materials and a very large fraction of molecular structures have
been described with it; Walter Kohn shared the 1998 Nobel Prize in Chemistry
for it; and it is now standard in quantum chemistry, materials science,
condensed-matter physics and, increasingly, low-energy nuclear physics.  It is
also, as we shall see, an approximation whose errors are of a quite different
character from those of the wave-function methods, and understanding that
difference is the point of putting it here.

%----------------------------------------------------------------------
\section{The electronic Hamiltonian}
\label{sec:electronichamiltonian}

\index{Born-Oppenheimer approximation}
Any piece of matter -- a crystal, a molecule, a living cell -- is a collection
of nuclei and electrons held together by the Coulomb interaction, and is
governed by
\begin{equation}
  \imath\hbar\frac{\partial}{\partial t}\Phi(\bm{r};t)
   = \left(-\sum_{i}^{N}\frac{\hbar^2}{2m_i}
     \frac{\partial^2}{\partial\bm{r}_i^2}
     + \sum_{i<j}^{N}\frac{e^2Z_iZ_j}{\vert\bm{r}_i-\bm{r}_j\vert}
     \right)\Phi(\bm{r};t).
  \label{eq:8-manybody}
\end{equation}
Nothing else is needed in principle, and nothing is soluble in practice.

The first simplification is universal.  Even the lightest nucleus is nearly
two thousand times heavier than an electron, so on the time scale of electronic
motion the nuclei may be treated as fixed.  This is the
\emph{Born-Oppenheimer approximation}, and it turns
Eq.~(\ref{eq:8-manybody}) into a time-independent problem for the electrons
alone, moving in the static field of the nuclei,
\begin{equation}
  \left[\sum_{i}^{N}\left(-\frac{\hbar^2}{2m}
    \frac{\partial^2}{\partial\bm{r}_i^2}
    + v_{\mathrm{ext}}(\bm{r}_i)\right)
    + \sum_{i<j}^{N}v(r_{ij})\right]\Psi(\bm{r}) = E_0\Psi(\bm{r}),
  \label{eq:8-electronic}
\end{equation}
with $r_{ij}=\vert\bm{r}_i-\bm{r}_j\vert$ and the external potential supplied
by the nuclei,
\begin{equation}
  v_{\mathrm{ext}}(\bm{r}_i) = -\sum_{j}\frac{e^2Z_j}{
    \vert\bm{r}_i-\bm{R}_j\vert}.
  \label{eq:8-vext}
\end{equation}
We have written a generic two-body interaction $v(r_{ij})$ rather than the
bare Coulomb repulsion, because everything below applies unchanged to any
interaction depending only on the relative distance -- nuclear forces
included, where $v_{\mathrm{ext}}$ is absent and the system is self-bound, as
Section~\ref{sec:whyhf} noted.

Equation~(\ref{eq:8-electronic}) is the Hamiltonian of
Eq.~(\ref{eq:4-generalH}) written in coordinate space.  Only one thing about
it will matter in what follows: the kinetic energy and the interaction are the
\emph{same for every system}, and the only thing that distinguishes one
material from another is $v_{\mathrm{ext}}$.

%----------------------------------------------------------------------
\section{Reduced density matrices}
\label{sec:reduceddensity}

\index{density matrix!one-body}\index{density matrix!two-body}
The central object of this chapter is the electronic density.  It is the
diagonal of the one-body reduced density matrix, which we met in
Section~\ref{sec:naturalorbitals} as a matrix in a single-particle basis and
now write in coordinate space.

\subsection{The one-body density}
\label{sec:onebodydensity}

Define the density operator
\begin{equation}
  \hat{n}(\bm{r}) = \sum_{i=1}^{N}\delta(\bm{r}-\bm{r}_i),
  \qquad
  n(\bm{r}) = \frac{\bracket{\Psi}{\hat{n}(\bm{r})}{\Psi}}
    {\braket{\Psi}{\Psi}},
  \label{eq:8-densityoperator}
\end{equation}
for any normalisable state $\Psi$.  Equivalently, integrating out all
coordinates but one,
\begin{equation}
  n(\bm{r}) = N\int \ud{^3\bm{r}_2}\cdots\ud{^3\bm{r}_N}\,
    \left\vert\Psi(\bm{r},\bm{r}_2,\dots,\bm{r}_N)\right\vert^2 .
  \label{eq:8-onebodydensity}
\end{equation}
For a single Slater determinant built from orbitals $\psi_i$ this collapses to
\begin{equation}
  n(\bm{r}) = \sum_{i=1}^{N}\left\vert\psi_i(\bm{r})\right\vert^2,
  \label{eq:8-densitydeterminant}
\end{equation}
a result already used in Section~\ref{sec:scf}.

In second quantization the operator reads
\begin{equation}
  \hat{n}(\bm{r}) = \sum_{\alpha\beta}
    \psi^{*}_{\alpha}(\bm{r})\psi_{\beta}(\bm{r})\,
    \hat{a}^{\dagger}_{\alpha}\hat{a}_{\beta},
  \label{eq:8-densitysecondquant}
\end{equation}
so that its expectation value in \emph{any} state -- a full
configuration-interaction ground state
$\bket{\Psi_0}=\sum_i C_{0i}\bket{\Phi_i}$ included -- is
\begin{equation}
  n(\bm{r}) = \sum_{ij}C^{*}_{0i}C_{0j}
    \sum_{\alpha\beta}\psi^{*}_{\alpha}(\bm{r})\psi_{\beta}(\bm{r})
    \bracket{\Phi_i}{\hat{a}^{\dagger}_{\alpha}\hat{a}_{\beta}}{\Phi_j}.
  \label{eq:8-densityfci}
\end{equation}
Integrating over $\bm{r}$ and using orthonormality gives
$\int n(\bm{r})\ud{^3\bm{r}}=N$, as it must.

The more general object, needed below, is the one-body reduced density matrix
off the diagonal,
\begin{equation}
  \rho^{(1)}(x';x) = N\int \ud{x_2}\cdots\ud{x_N}\,
    \Psi^{*}(x',x_2,\dots,x_N)\,\Psi(x,x_2,\dots,x_N),
  \label{eq:8-rho1}
\end{equation}
with $x$ the combined space and spin coordinate.  It is Hermitian, its
diagonal is the density, $n(x)=\rho^{(1)}(x;x)$, and its eigenvalues are the
occupation numbers of Section~\ref{sec:naturalorbitals}: at most one for
fermions, and exactly $N$ ones and the rest zeros when -- and only when --
$\Psi$ is a single Slater determinant, in which case $\rho^{(1)}$ is the
projector onto the occupied subspace.

\subsection{The two-body density}
\label{sec:twobodydensity}

Two-body operators need the pair density,
\begin{equation}
  \rho^{(2)}(x_1,x_2;y_1,y_2) = \frac{N(N-1)}{2}
    \int \ud{x_3}\cdots\ud{x_N}\,
    \Psi^{*}(y_1,y_2,x_3,\dots,x_N)\,\Psi(x_1,x_2,x_3,\dots,x_N),
  \label{eq:8-rho2}
\end{equation}
the prefactor counting the $N(N-1)/2$ distinct pairs.  Its diagonal
$\rho^{(2)}(x_1,x_2)$ is the joint probability of finding one particle at
$x_1$ and another at $x_2$, and it contains everything needed for the
expectation value of any two-body operator -- which is to say, everything
needed for the energy.  In second quantization,
\begin{equation}
  \hat{\rho}^{(2)}(\bm{r}_1,\bm{r}_2)
   = \sum_{\alpha\beta\gamma\delta}
     \rho_{\alpha\gamma}(\bm{r}_1)\rho_{\beta\delta}(\bm{r}_2)\,
     \hat{a}^{\dagger}_{\alpha}\hat{a}^{\dagger}_{\beta}
     \hat{a}_{\delta}\hat{a}_{\gamma},
  \qquad
  \rho_{\alpha\beta}(\bm{r})
   = \psi^{*}_{\alpha}(\bm{r})\psi_{\beta}(\bm{r}).
  \label{eq:8-rho2secondquant}
\end{equation}

\subsection{Antisymmetry, and when the pair density factorises}
\label{sec:factorisation}

Fermionic wave functions change sign under the exchange of any two particles,
\begin{equation}
  \Psi(\dots,x_i,\dots,x_j,\dots)
   = -\Psi(\dots,x_j,\dots,x_i,\dots),
  \label{eq:8-antisymmetry}
\end{equation}
and this has a direct consequence for the pair density: it must vanish when
$x_1=x_2$, since two identical fermions cannot be at the same point in the
same spin state.  This is the Pauli principle expressed as a statement about
$\rho^{(2)}$, and any approximate treatment that violates it is in trouble.

For a single Slater determinant the pair density is completely determined by
the one-body density matrix.  Applying Wick's theorem,
Section~\ref{sec:wicktheorem}, to $\bracket{\Phi_0}{
\hat{a}^{\dagger}\hat{a}^{\dagger}\hat{a}\hat{a}}{\Phi_0}$ leaves exactly two
contractions, and
\begin{equation}
  \boxed{\;
  \rho^{(2)}(x_1,x_2;y_1,y_2)
   = \rho^{(1)}(x_1;y_1)\rho^{(1)}(x_2;y_2)
   - \rho^{(1)}(x_1;y_2)\rho^{(1)}(x_2;y_1).
  \;}
  \label{eq:8-factorisation}
\end{equation}
The first term is the naive product -- what one would write for independent
particles -- and the second is the exchange term that enforces antisymmetry.
On the diagonal,
\begin{equation}
  \rho^{(2)}(x_1,x_2) = n(x_1)n(x_2)
    - \left\vert\rho^{(1)}(x_1;x_2)\right\vert^2 ,
  \label{eq:8-pairdiagonal}
\end{equation}
and setting $x_1=x_2$ gives $n(x_1)^2-n(x_1)^2=0$ exactly, as required.

Equation~(\ref{eq:8-factorisation}) is worth dwelling on, because it marks
precisely the boundary of mean-field theory.  For a Hartree product -- a
plain product of orbitals with no antisymmetry -- one would have
$\rho^{(2)}=n(x_1)n(x_2)$, but such a state is not a legitimate fermionic wave
function.  The closest legitimate object is a single determinant, for which
Eq.~(\ref{eq:8-factorisation}) holds \emph{exactly}, with the exchange term as
the only two-particle correlation.  Beyond that -- for any genuinely
correlated state -- the pair density cannot be reduced to one-body quantities
at all, and the difference is what we have been calling correlation
throughout this book.

\begin{notebox}
\textbf{Where this leaves us.}  The energy of Eq.~(\ref{eq:8-electronic})
needs only $\rho^{(1)}$ and $\rho^{(2)}$; the full wave function is
superfluous.  That observation is the origin of reduced-density-matrix
methods, which try to determine $\rho^{(2)}$ directly.  The difficulty is
the \emph{$N$-representability problem}: not every Hermitian, antisymmetric,
correctly normalised $\rho^{(2)}$ comes from an actual $N$-fermion wave
function, and characterising those that do is extremely hard.  Density
functional theory sidesteps this by going all the way down to $n(\bm{r})$, at
the price of an unknown functional.  The two difficulties -- an unknown
functional here, an unknown representability condition there -- are two faces
of the same problem, and neither has been solved.
\end{notebox}

%----------------------------------------------------------------------
\section{The Hohenberg-Kohn theorems}
\label{sec:hohenbergkohn}

\index{Hohenberg-Kohn theorems}
We now come to the two results on which everything rests.  Consider the family
of Hamiltonians
\begin{equation}
  \hat{H} = \hat{T} + \hat{V}_{\mathrm{ext}} + \hat{V},
  \qquad
  \hat{V}_{\mathrm{ext}}\in\mathcal{V}_{\mathrm{ext}},
  \label{eq:8-family}
\end{equation}
in which $\hat{T}$ and $\hat{V}$ are fixed once and for all -- they are
properties of electrons, not of materials -- and only the external potential
varies.  Assume for now that the ground state is non-degenerate.

\begin{theorem}[Hohenberg-Kohn I]
The external potential $v_{\mathrm{ext}}(\bm{r})$ is determined, up to an
additive constant, by the ground-state density $n_0(\bm{r})$.
\end{theorem}

\begin{theorem}[Hohenberg-Kohn II]
There exists a universal functional $F[n]$, the same for every external
potential, such that
\begin{equation}
  E[n] = F[n] + \int \ud{^3\bm{r}}\,v_{\mathrm{ext}}(\bm{r})n(\bm{r}),
  \qquad
  F[n] = T[n] + V_{\mathrm{int}}[n],
  \label{eq:8-hkfunctional}
\end{equation}
is minimised, over densities integrating to $N$, by the true ground-state
density, and its minimum value is the ground-state energy $E_0$.
\end{theorem}

Together these say something startling.  The density -- three variables --
determines the potential, hence the Hamiltonian, hence the full many-body
wave function and every ground-state observable.  A quantity that is measured
routinely by X-ray diffraction contains, in principle, all the information in
the $3N$-dimensional $\Psi$.

\subsection{Proof of the first theorem}
\label{sec:hkproof}

The proof is a two-line \emph{reductio ad absurdum} and is worth giving in
full, because its simplicity is itself informative.

Suppose two external potentials $v^{(1)}$ and $v^{(2)}$, differing by more
than a constant, produced the same ground-state density $n(\bm{r})$.  They have
different Hamiltonians $\hat{H}^{(1)},\hat{H}^{(2)}$ and, by the unique
continuation theorem, different ground states
$\Psi^{(1)}\ne\Psi^{(2)}$ with energies $E^{(1)},E^{(2)}$.

Use $\Psi^{(2)}$ as a trial function for $\hat{H}^{(1)}$.  The variational
principle of Section~\ref{sec:variational} is strict for a non-degenerate
ground state, so
\begin{equation}
  E^{(1)} < \bracket{\Psi^{(2)}}{\hat{H}^{(1)}}{\Psi^{(2)}}
   = \bracket{\Psi^{(2)}}{\hat{H}^{(2)}}{\Psi^{(2)}}
   + \int\left[v^{(1)}-v^{(2)}\right]n(\bm{r})\ud{^3\bm{r}}
   = E^{(2)} + \int\left[v^{(1)}-v^{(2)}\right]n .
  \label{eq:8-hkineq1}
\end{equation}
The two Hamiltonians differ only in the external potential, and the density is
by hypothesis the same, which is what allows the middle step.  Interchanging
the labels,
\begin{equation}
  E^{(2)} < E^{(1)} + \int\left[v^{(2)}-v^{(1)}\right]n .
  \label{eq:8-hkineq2}
\end{equation}
Adding Eqs.~(\ref{eq:8-hkineq1}) and~(\ref{eq:8-hkineq2}), the integrals
cancel and we are left with
\begin{equation}
  E^{(1)} + E^{(2)} < E^{(1)} + E^{(2)},
  \label{eq:8-hkcontradiction}
\end{equation}
which is absurd.  Hence no two potentials differing by more than a constant
can share a ground-state density.

The second theorem follows immediately.  Since $n_0$ fixes
$v_{\mathrm{ext}}$, it fixes $\hat{H}$ and therefore $\Psi[n]$, so
\begin{equation}
  F_{\mathrm{HK}}[n] \equiv \bracket{\Psi[n]}{\hat{T}+\hat{V}}{\Psi[n]}
  \label{eq:8-fhk}
\end{equation}
is a well-defined functional of the density alone, and it is universal
because $\hat{T}$ and $\hat{V}$ contain no reference to
$v_{\mathrm{ext}}$.  The Rayleigh-Ritz principle applied within this family
gives $E_0=\min_{n}E[n]$.

\subsection{What the theorems do not give}
\label{sec:hkcaveats}

Three caveats deserve stating plainly, because they are the difference
between a beautiful theorem and a working method.

\paragraph{They are existence theorems.}
Nothing in the proof constructs $F[n]$, and nothing suggests it is simple.
All the many-body difficulty of the preceding chapters is still there; it has
been moved inside $F[n]$.  This is why the practical history of density
functional theory is a history of \emph{approximate functionals}.

\paragraph{$v$-representability.}
The argument assumes that the density under consideration is the ground-state
density of \emph{some} potential in the family.  Not every well-behaved,
positive, correctly normalised $n(\bm{r})$ is, and characterising those that
are is unsolved.  Levy and Lieb removed the difficulty by redefining the
functional as a constrained search over wave functions,
\begin{equation}
  F_{\mathrm{LL}}[n] = \min_{\Psi\to n}
    \bracket{\Psi}{\hat{T}+\hat{V}}{\Psi},
  \label{eq:8-levylieb}
\end{equation}
the minimum running over all antisymmetric $\Psi$ that yield the density $n$.
This is defined for any $N$-representable density, agrees with
$F_{\mathrm{HK}}$ where the latter exists, and extends the theory to
degenerate ground states.  It is also, of course, no more computable.

\paragraph{Ground states only.}
The functional delivers $E_0$ and $n_0$ and nothing else.  Excitation
energies -- the subject of Chapter~\ref{chap:mfapplications} -- are not
obtained from the ground-state functional, and reaching them requires a
separate construction such as time-dependent density functional theory.  The
Kohn-Sham eigenvalues introduced below are \emph{not} excitation energies, a
point we return to in Section~\ref{sec:ksmeaning}.

\subsection{The variational equation}
\label{sec:dftvariational}

Given $F[n]$, the ground state follows from minimising $E[n]$ subject to
$\int n = N$.  Introducing a Lagrange multiplier exactly as in
Section~\ref{sec:varcalculus},
\begin{equation}
  \frac{\delta}{\delta n(\bm{r})}
  \left\{E[n]-\mu\left(\int n\,\ud{^3\bm{r}}-N\right)\right\} = 0
  \qquad\Longrightarrow\qquad
  \left.\frac{\delta E[n]}{\delta n(\bm{r})}\right\vert_{n=n_0} = \mu ,
  \label{eq:8-eulerequation}
\end{equation}
so that the functional derivative of the energy is constant throughout the
system.  That constant is the chemical potential, the Fermi level.  The
pattern is by now familiar: a constrained variation produces a stationarity
condition, and the multiplier enforcing normalisation is the energy conjugate
to the particle number.

%----------------------------------------------------------------------
\section{Building a functional from the electron gas}
\label{sec:functionalfromgas}

\index{Thomas-Fermi model}\index{local density approximation}
The theorems tell us $F[n]$ exists.  To make progress we need an
approximation, and the oldest and still the most instructive strategy is to
take a system we \emph{can} solve exactly and use it locally.

That system is the homogeneous electron gas of
Section~\ref{sec:electrongas}.  There we computed, for a uniform density
$n$, both the kinetic energy per unit volume,
\begin{equation}
  t_0[n] = \frac{3\hbar^2(3\pi^2)^{2/3}}{10m}\,n^{5/3},
  \label{eq:8-tf}
\end{equation}
which is Eq.~(\ref{eq:6-egkineticdensity}) divided by the volume, and the
exchange energy per particle.  If the density varies slowly enough, we may
pretend that each small region behaves like a uniform gas at the local
density and integrate,
\begin{equation}
  \mathcal{T}_0[n] = \frac{3\hbar^2(3\pi^2)^{2/3}}{10m}
    \int \ud{^3\bm{r}}\,\left(n(\bm{r})\right)^{5/3}.
  \label{eq:8-tffunctional}
\end{equation}
This is the \emph{Thomas-Fermi} kinetic functional, and with the classical
electrostatic (Hartree) energy
\begin{equation}
  \mathcal{U}_{\mathrm{H}}[n] = \frac{e^2}{2}
    \iint \ud{^3\bm{r}}\ud{^3\bm{r}'}\,
    \frac{n(\bm{r})n(\bm{r}')}{\vert\bm{r}-\bm{r}'\vert},
  \label{eq:8-hartree}
\end{equation}
it gives a complete, if crude, theory.  Everything not accounted for is
collected into a single remainder, the \emph{exchange-correlation energy}
$E_{\mathrm{xc}}[n]$, and the Hohenberg-Kohn functional is written
\begin{equation}
  \mathcal{E}_{\mathrm{HK}}[n]
   = \mathcal{T}_0[n]
   + \int \ud{^3\bm{r}}\,v_{\mathrm{ext}}(\bm{r})n(\bm{r})
   + \mathcal{U}_{\mathrm{H}}[n]
   + E_{\mathrm{xc}}[n].
  \label{eq:8-hkdecomposition}
\end{equation}
The decomposition is a definition, not an approximation: whatever is missing
is in $E_{\mathrm{xc}}$ by construction.  It is useful only if that remainder
is small, which is the whole gamble of the subject.

\paragraph{Why Thomas-Fermi is not enough.}
Equation~(\ref{eq:8-tffunctional}) has no orbitals in it, and minimising
Eq.~(\ref{eq:8-hkdecomposition}) directly over $n$ gives an algebraic
equation.  It is enormously cheaper than anything in the preceding chapters.
It also fails, and the failure is instructive rather than marginal.  The
program \texttt{dft.py} compares the exact kinetic energy of $N$
non-interacting fermions in a one-dimensional trap with the local estimate:

\begin{center}
\renewcommand{\arraystretch}{1.2}
\setlength{\tabcolsep}{11pt}
\begin{tabular}{rrrr}
\toprule
$N$ & $T$ (exact) & $T$ (Thomas-Fermi) & ratio\\
\midrule
$1$ & $0.250000$  & $0.302300$  & $1.2092$\\
$2$ & $1.000000$  & $1.074844$  & $1.0748$\\
$3$ & $2.250000$  & $2.340025$  & $1.0400$\\
$4$ & $4.000000$  & $4.101805$  & $1.0255$\\
$6$ & $9.000000$  & $9.120038$  & $1.0133$\\
$8$ & $16.000000$ & $16.134286$ & $1.0084$\\
\bottomrule
\end{tabular}
\end{center}

The error falls steadily as the particle number grows, which is exactly what
the construction promises -- the density becomes smoother on the scale of the
local Fermi wavelength -- but a one per cent error in the \emph{largest} term
of the energy is fatal when chemical accuracy is wanted.  Worse, the
Thomas-Fermi density has no shell structure at all: it is smooth and
featureless, with none of the nodal structure that orbitals supply, and
Teller proved that within the model no molecule is bound.  Something has to
be done about the kinetic energy specifically.

%----------------------------------------------------------------------
\section{The Kohn-Sham construction}
\label{sec:kohnsham}

\index{Kohn-Sham equations}
Kohn and Sham's idea was to keep orbitals for the one term that cannot be
approximated locally, and only that term.

Introduce an auxiliary system of $N$ \emph{non-interacting} fermions moving
in some effective potential $v_{\mathrm{KS}}(\bm{r})$, chosen so that its
ground-state density equals the density of the real, interacting system.  Its
kinetic energy,
\begin{equation}
  T_s[n] = -\frac{\hbar^2}{2m}\sum_{i=1}^{N}
    \int \ud{^3\bm{r}}\,\psi^{*}_i(\bm{r})\nabla^2\psi_i(\bm{r}),
  \label{eq:8-Ts}
\end{equation}
is computed exactly from orbitals rather than approximated from the density,
and it captures the shell structure that Thomas-Fermi threw away.  It is not
the true kinetic energy $T[n]$ of the interacting system, but the difference
is small, and we simply redefine the remainder to absorb it,
\begin{equation}
  E_{\mathrm{xc}}[n] \equiv \left(T[n]-T_s[n]\right)
    + \left(V_{\mathrm{int}}[n]-\mathcal{U}_{\mathrm{H}}[n]\right).
  \label{eq:8-excdefinition}
\end{equation}
The total energy is then
\begin{equation}
  E[n] = T_s[n]
    + \int \ud{^3\bm{r}}\,v_{\mathrm{ext}}n
    + \mathcal{U}_{\mathrm{H}}[n]
    + E_{\mathrm{xc}}[n].
  \label{eq:8-ksenergy}
\end{equation}

Varying Eq.~(\ref{eq:8-ksenergy}) with respect to the orbitals, subject to
orthonormality, and following exactly the steps of
Section~\ref{sec:hfcoefficients}, gives the \emph{Kohn-Sham equations}
\begin{equation}
  \boxed{\;
  \left(-\frac{\hbar^2}{2m}\nabla^2
    + v_{\mathrm{KS}}(\bm{r})\right)\psi_i(\bm{r})
   = \varepsilon_i\psi_i(\bm{r}),
  \;}
  \label{eq:8-kohnsham}
\end{equation}
with the Kohn-Sham potential
\begin{equation}
  v_{\mathrm{KS}}(\bm{r}) = v_{\mathrm{ext}}(\bm{r})
    + \int \ud{^3\bm{r}'}\,\frac{e^2 n(\bm{r}')}{\vert\bm{r}-\bm{r}'\vert}
    + v_{\mathrm{xc}}(\bm{r}),
  \qquad
  v_{\mathrm{xc}}(\bm{r})
   = \frac{\delta E_{\mathrm{xc}}[n]}{\delta n(\bm{r})} ,
  \label{eq:8-kspotential}
\end{equation}
and the density reconstructed from the occupied orbitals,
\begin{equation}
  n(\bm{r}) = \sum_{i=1}^{N}\vert\psi_i(\bm{r})\vert^2 .
  \label{eq:8-ksdensity}
\end{equation}
Since $v_{\mathrm{KS}}$ depends on $n$ and $n$ on the orbitals, the equations
are nonlinear and must be iterated to self-consistency -- the same
self-consistent field loop as Section~\ref{sec:scf}, and for the same reason.

\paragraph{The total energy.}
The sum of the Kohn-Sham eigenvalues is \emph{not} the total energy, exactly
as in Hartree-Fock, Exercise~\ref{ex:6-hfequations}b).  Since
$\sum_i\varepsilon_i = T_s + \int v_{\mathrm{KS}}n$,
\begin{equation}
  E = \sum_{i=1}^{N}\varepsilon_i
    - \mathcal{U}_{\mathrm{H}}[n]
    + E_{\mathrm{xc}}[n]
    - \int \ud{^3\bm{r}}\,v_{\mathrm{xc}}(\bm{r})n(\bm{r}) ,
  \label{eq:8-kstotalenergy}
\end{equation}
the last three terms correcting the double counting of the Hartree and
exchange-correlation contributions.

\subsection{Kohn-Sham against Hartree-Fock}
\label{sec:ksvshf}

It is worth setting the two schemes side by side, since they look almost
identical and differ in one decisive respect.  The Hartree-Fock equation of
Chapter~\ref{chap:hartrefock} reads, in coordinate space,
\begin{align}
  \Bigg(-\frac{\hbar^2}{2m}\nabla^2 + v_{\mathrm{ext}}(\bm{r})
   &+ \int \ud{^3\bm{r}'}\,w(\bm{r},\bm{r}')n(\bm{r}')\Bigg)\phi_k(\bm{r})
   \nonumber\\
   &\underbrace{-\sum_{l=1}^{N}\int \ud{^3\bm{r}'}\,
     \phi^{*}_l(\bm{r}')w(\bm{r},\bm{r}')\phi_k(\bm{r}')\phi_l(\bm{r})}
     _{\text{exchange}}
   = \varepsilon_k\phi_k(\bm{r}),
  \label{eq:8-hfcoordinate}
\end{align}
whereas the Kohn-Sham equation is
\begin{equation}
  \left(-\frac{\hbar^2}{2m}\nabla^2 + v_{\mathrm{ext}}(\bm{r})
    + \int \ud{^3\bm{r}'}\,w(\bm{r},\bm{r}')n(\bm{r}')
    + \underbrace{v_{\mathrm{xc}}([n];\bm{r})}
      _{\text{exchange {\it and} correlation}}\right)\phi_k(\bm{r})
   = \varepsilon_k\phi_k(\bm{r}).
  \label{eq:8-kscoordinate}
\end{equation}
The difference is the last term.  Hartree-Fock's exchange is \emph{non-local}
-- an integral operator, different for each orbital -- and it is exact.
Kohn-Sham's $v_{\mathrm{xc}}$ is a \emph{local multiplicative potential}, the
same for every orbital, and it is approximate; but it also contains
correlation, which Hartree-Fock does not have at all.  The trade is
non-locality and exactness against locality, cheapness and a bit of
correlation.  Both scale formally as the cube of the number of basis
functions, but the Kohn-Sham prefactor is far smaller because no
four-index integrals are needed for $v_{\mathrm{xc}}$.

\begin{notebox}
\textbf{What the Kohn-Sham orbitals are.}
\label{sec:ksmeaning}
They are the orbitals of a fictitious non-interacting system chosen to
reproduce the right density.  They are not approximations to the true
one-electron wave functions, and the eigenvalues $\varepsilon_i$ are not
ionisation energies -- Koopmans' theorem, Exercise~\ref{ex:6-hfequations}c),
applies to Hartree-Fock, not here.  The one exception is the highest occupied
level, which for the exact functional equals minus the exact ionisation
energy.  The systematic underestimation of band gaps by Kohn-Sham eigenvalue
differences is not a numerical deficiency but a consequence of asking the
orbitals a question they were not built to answer.
\end{notebox}

%----------------------------------------------------------------------
\section{The local density approximation}
\label{sec:lda}

\index{local density approximation}\index{exchange-correlation energy}
Everything now rests on $E_{\mathrm{xc}}[n]$.  The oldest approximation, and
still the reference point for all the others, is the \emph{local density
approximation}: assume that the exchange-correlation energy density at a
point depends only on the density at that point, and take that dependence
from the uniform electron gas,
\begin{equation}
  E_{\mathrm{xc}}^{\mathrm{LDA}}[n]
   = \int \ud{^3\bm{r}}\,n(\bm{r})\,
     \epsilon_{\mathrm{xc}}\!\left(n(\bm{r})\right),
  \label{eq:8-lda}
\end{equation}
where $\epsilon_{\mathrm{xc}}(n)$ is the exchange-correlation energy per
particle of a uniform gas of density $n$.  The corresponding potential
follows by differentiation,
\begin{equation}
  v_{\mathrm{xc}}(\bm{r})
   = \frac{\ud{}}{\ud{n}}
     \left[n\,\epsilon_{\mathrm{xc}}(n)\right]_{n=n(\bm{r})} .
  \label{eq:8-vxclda}
\end{equation}

The exchange part is available in closed form.  It is the exchange energy of
Section~\ref{sec:egsingleparticle} expressed per particle rather than per
state,
\begin{equation}
  \epsilon_{x}(n) = -\frac{3}{4}\left(\frac{3}{\pi}\right)^{1/3}n^{1/3},
  \qquad
  v_{x}(n) = \frac{4}{3}\epsilon_x(n),
  \label{eq:8-ldaexchange}
\end{equation}
in Hartree atomic units.  That this must reproduce the electron-gas result of
Chapter~\ref{chap:hartrefock} is a consistency check the program performs:

\begin{center}
\renewcommand{\arraystretch}{1.2}
\setlength{\tabcolsep}{11pt}
\begin{tabular}{rrrr}
\toprule
$r_s$ & $\epsilon_x$ [Ha] & $\epsilon_x$ [Ry] & $-0.916/r_s$ [Ry]\\
\midrule
$1.0000$ & $-0.458165$ & $-0.916331$ & $-0.916000$\\
$2.0000$ & $-0.229083$ & $-0.458165$ & $-0.458000$\\
$4.0000$ & $-0.114541$ & $-0.229083$ & $-0.229000$\\
$4.8253$ & $-0.094951$ & $-0.189901$ & $-0.189833$\\
$6.0000$ & $-0.076361$ & $-0.152722$ & $-0.152667$\\
\bottomrule
\end{tabular}
\end{center}

The coefficient comes out as $0.916331$ against the $0.9163$ of
Eq.~(\ref{eq:6-egcoefficients}).  This agreement is a tautology -- the local
approximation is \emph{defined} to be exact for the uniform gas -- and the
interesting question is not whether it holds but why the approximation
remains useful when the density is far from uniform.

The correlation part $\epsilon_c(n)$ has no closed form.  It is obtained from
quantum Monte Carlo calculations of the uniform gas, the diffusion Monte
Carlo results of Ceperley and Alder being the standard source, and then
fitted to a convenient analytic form.  This is worth noticing: the input to
the most widely used approximation in computational materials science is a
wave-function calculation of the kind developed in the preceding chapters.

\subsection{The exchange hole}
\label{sec:exchangehole}

\index{exchange hole}
The reason a local approximation works at all becomes clearer if one looks at
what exchange actually does.  Divide the pair
density~(\ref{eq:8-pairdiagonal}) by the product of densities to define the
pair correlation function $g$; for the uniform gas at the Hartree-Fock level
it depends only on $y=k_F s$ and is known exactly,
\begin{equation}
  g(y) = 1 - \frac{9}{2}
    \left[\frac{\sin y - y\cos y}{y^3}\right]^2 .
  \label{eq:8-paircorrelation}
\end{equation}
The program tabulates it:

\begin{center}
\renewcommand{\arraystretch}{1.2}
\setlength{\tabcolsep}{11pt}
\begin{tabular}{rll}
\toprule
$k_F s$ & $g(s)$ & \\
\midrule
$0.00$  & $0.500000$ & same spins never coincide\\
$0.50$  & $0.524471$ & \\
$1.00$  & $0.591838$ & \\
$2.00$  & $0.786732$ & \\
$4.00$  & $0.996208$ & \\
$8.00$  & $0.999920$ & the hole has healed\\
\bottomrule
\end{tabular}
\end{center}

At coincidence $g(0)=\tfrac12$: electrons of the same spin are completely
excluded, those of opposite spin are uncorrelated at this level.  Each
electron therefore moves inside a hole in the density of its neighbours, and
that hole obeys an exact sum rule -- it contains precisely one missing
electron,
\begin{equation}
  n\int \ud{^3\bm{s}}\left[g(s)-1\right]
   = \frac{4}{3\pi}\int_0^{\infty}\ud{y}\,y^2\left[g(y)-1\right]
   = -1 ,
  \label{eq:8-sumrule}
\end{equation}
which the program verifies numerically to four digits.

This is the deeper reason the local approximation succeeds.  The
exchange-correlation energy is, to a good approximation, the electrostatic
interaction of each electron with its own hole.  That energy depends mostly
on the \emph{total charge} of the hole -- which the sum rule fixes exactly --
and only weakly on its shape.  A functional that gets the shape badly wrong
but the normalisation right can still give good energies, and this is what
the local density approximation does.

\subsection{Beyond LDA}
\label{sec:beyondlda}

The density in a real atom or molecule is anything but slowly varying, and
the natural first correction is to let the functional see the gradient as
well,
\begin{equation}
  E_{\mathrm{xc}}^{\mathrm{GGA}}[n]
   = \int \ud{^3\bm{r}}\,n(\bm{r})\,
     \epsilon_{\mathrm{xc}}\!\left(n(\bm{r}),
       \vert\nabla n(\bm{r})\vert\right),
  \label{eq:8-gga}
\end{equation}
the \emph{generalised gradient approximation}.  A naive gradient expansion of
the uniform gas actually makes things worse; what works is to construct the
gradient dependence so as to preserve the exact constraints -- the sum rule
above, scaling relations, known limits -- and this is how the widely used
functionals were built.  Beyond that lie meta-GGAs, which add the kinetic
energy density, and hybrids, which mix in a fraction of exact Hartree-Fock
exchange, Eq.~(\ref{eq:8-hfcoordinate}), to repair the self-interaction error
discussed next.

%----------------------------------------------------------------------
\section{Kohn-Sham in practice}
\label{sec:kspractice}

To see the machinery work, and fail, we apply it to a system we have solved
twice already: $N$ spinless fermions in a one-dimensional harmonic trap with
a softened Coulomb repulsion, in a basis of eight oscillator orbitals.  This
is the system of Section~\ref{sec:energyfunctional} and of the Hartree-Fock
calculation of Section~\ref{sec:scf}, so the three treatments can be compared
for the same Hamiltonian.

There is no tabulated exchange functional for a one-dimensional gas with this
interaction, so the program builds one, which is instructive in itself.  For a
spinless one-dimensional gas the Fermi momentum is $k_F=\pi n$ and the
one-body density matrix is $\rho^{(1)}(s)=\sin(k_Fs)/\pi s$, so
\begin{equation}
  \epsilon_x(n) = -\frac{1}{n}\int_0^{\infty}\ud{s}\,v(s)
    \left[\frac{\sin(k_F s)}{\pi s}\right]^2 ,
  \label{eq:8-lda1d}
\end{equation}
which is evaluated numerically on a grid of densities, tabulated and
interpolated.  This is precisely how real functionals are made.  The result:

\begin{center}
\renewcommand{\arraystretch}{1.2}
\setlength{\tabcolsep}{11pt}
\begin{tabular}{rrr}
\toprule
$n$ & $\epsilon_x(n)$ & $v_x(n)$\\
\midrule
$0.05$ & $-0.17352608$ & $-0.29740316$\\
$0.10$ & $-0.27855522$ & $-0.45906496$\\
$0.20$ & $-0.42332628$ & $-0.65789106$\\
$0.40$ & $-0.59557914$ & $-0.84878595$\\
$0.80$ & $-0.75884371$ & $-0.96587633$\\
$1.50$ & $-0.86548980$ & $-0.99703512$\\
$2.50$ & $-0.91895548$ & $-0.99989688$\\
\bottomrule
\end{tabular}
\end{center}

Solving Eq.~(\ref{eq:8-kohnsham}) self-consistently in the oscillator basis
then gives

\begin{table}[h]
\centering
\renewcommand{\arraystretch}{1.25}
\setlength{\tabcolsep}{11pt}
\begin{tabular}{rrrr}
\toprule
$N$ & $E$ (Kohn-Sham) & $E$ (Hartree-Fock) & $E$ (exact)\\
\midrule
$1$ & $0.54752674$  & $0.50000000$  & $0.50000000$\\
$2$ & $2.72022768$  & $2.66308878$  & $2.65923049$\\
$3$ & $6.45108052$  & $6.39016133$  & ---\\
$4$ & $11.68602229$ & $11.62305385$ & ---\\
\bottomrule
\end{tabular}
\caption{Kohn-Sham with a local exchange functional, Hartree-Fock and the
exact answer for $N$ spinless fermions in a one-dimensional trap with a
softened Coulomb repulsion, in a basis of eight oscillator orbitals.  The
Hartree-Fock column reproduces Table~\ref{tab:hfgain}; the exact value for
$N=2$ comes from full diagonalisation in the same basis.}
\label{tab:ksresults}
\end{table}

Note first that the exact two-particle energy, $2.65923049$, lies
$-3.86\times10^{-3}$ below the Hartree-Fock value, which is the correlation
energy quoted in Eq.~(\ref{eq:2-correlation}) -- the three chapters are
consistent.  Note next that Kohn-Sham lies \emph{above} both.  This is not a
failure of the variational principle, because we have used exchange only and
put no correlation functional in at all; it is the error of treating exchange
locally, compounded by a second defect that we now isolate.

\subsection{The self-interaction error}
\label{sec:selfinteraction}

\index{self-interaction error}
Take a single particle.  There is nothing for it to interact with, so the
exact energy is the oscillator ground state, $0.5$.  Hartree-Fock gives
exactly that: in Eq.~(\ref{eq:8-hfcoordinate}) the direct term with $l=k$ is
cancelled term by term by the exchange term with $l=k$, and a Slater
determinant is rigorously free of self-interaction.

A local functional cannot do this, because it sees only the density, and the
density of one electron is indistinguishable from the density of half of two.
The program reports, for $N=1$:
\begin{equation}
  \mathcal{U}_{\mathrm{H}} = 0.60855522,
  \qquad
  E_x^{\mathrm{LDA}} = -0.56166291,
  \qquad
  \mathcal{U}_{\mathrm{H}} + E_x^{\mathrm{LDA}} = 0.04689232,
  \label{eq:8-sie}
\end{equation}
where the sum should vanish identically.  The residue $0.047$ is the
\emph{self-interaction error}, and it accounts for almost all of the
$0.0475$ by which Kohn-Sham misses the exact one-particle energy.

This single number explains a long list of known LDA failures: anions are
overbound, reaction barriers too low, band gaps too small, and charge is
over-delocalised, because an electron is spuriously repelled by itself and
spreads out to reduce that repulsion.  Self-interaction corrections remove it
explicitly, orbital by orbital; hybrid functionals reduce it by mixing in a
fraction of exact exchange, which is self-interaction free by construction.

\subsection{How local is too local?}
\label{sec:howlocal}

We can isolate the error of \emph{locality} alone, cleanly, by evaluating on
one and the same density both the exact exchange energy of the determinant
and the local estimate:

\begin{center}
\renewcommand{\arraystretch}{1.2}
\setlength{\tabcolsep}{11pt}
\begin{tabular}{rrrr}
\toprule
$N$ & $E_x$ (exact) & $E_x$ (local) & ratio\\
\midrule
$1$ & $-0.60855522$ & $-0.56166291$ & $0.9229$\\
$2$ & $-1.33077771$ & $-1.27749456$ & $0.9600$\\
$3$ & $-2.11087395$ & $-2.05548416$ & $0.9738$\\
$4$ & $-2.92523775$ & $-2.86866697$ & $0.9807$\\
\bottomrule
\end{tabular}
\end{center}

The local approximation recovers $92\%$ of the exchange energy for one
particle and $98\%$ for four, improving steadily as the density smooths out
on the scale of the local Fermi wavelength.  This is the behaviour the
construction promises, and the corresponding figure for real atoms in three
dimensions is close to $90\%$ -- an error of the order of ten per cent in a
large term, which is exactly why gradient corrections were needed before
density functional theory became quantitative for chemistry.

%----------------------------------------------------------------------
\section{Summary}
\label{sec:ch8summary}

Density functional theory rests on an exact statement -- the ground-state
density determines the Hamiltonian -- and on an approximation that the theorem
itself does not supply.  The Hohenberg-Kohn theorems guarantee a universal
functional $F[n]$ and a variational principle for it, but they are existence
theorems; all the many-body difficulty of the preceding chapters survives,
relocated inside $F[n]$.

The Kohn-Sham construction is what makes the theory workable.  By
reintroducing orbitals for the kinetic energy -- the one contribution that a
local approximation handles badly, as Thomas-Fermi showed -- it keeps shell
structure exactly and leaves only a small remainder to be approximated.  What
is left is a self-consistent one-body problem of exactly the same form as
Hartree-Fock, Eq.~(\ref{eq:6-hfequations}), differing in a single term: a
local multiplicative $v_{\mathrm{xc}}$ in place of the non-local exchange
operator.

The local density approximation takes that remainder from the electron gas of
Chapter~\ref{chap:hartrefock}, and its success rests less on the density
being slowly varying -- it never is -- than on the exchange-correlation hole
sum rule, which fixes the quantity the energy is most sensitive to.  Its two
characteristic failures are visible in a system with four particles: exchange
treated locally is wrong by a few per cent, and the self-interaction that
Hartree-Fock cancels identically survives.

It is worth being clear about where this method sits relative to the others in
this book.  Configuration interaction, coupled cluster and the random-phase
approximation are systematically improvable: there is a hierarchy, and one
can in principle climb it towards the exact answer.  Density functional
theory is not of that kind.  There is no sequence of functionals converging to
$F[n]$, and no internal estimate of the error.  What it offers instead is a
cost that scales gently enough to treat thousands of atoms, and an accuracy
that -- for reasons that are now partly understood -- is far better than the
crudeness of the approximations would suggest.

The program \texttt{dft.py} performs every calculation of this chapter, and
the notebook \texttt{dft.ipynb} runs it cell by cell.

%----------------------------------------------------------------------
\section{Exercises}
\label{sec:ch8exercises}
%=============================================================================

\subsection*{Exercise 1: densities and density matrices}
\label{ex:8-densities}
\addcontentsline{toc}{subsection}{Exercise 1: densities and density matrices}

\begin{enumerate}
  \item[a)] Show that Eq.~(\ref{eq:8-densityoperator}) and
  Eq.~(\ref{eq:8-onebodydensity}) define the same quantity, and that
  $\int n(\bm{r})\ud{^3\bm{r}}=N$.
  \item[b)] Derive Eq.~(\ref{eq:8-densitydeterminant}) for a single Slater
  determinant, and show that $\rho^{(1)}$ is then the projector onto the
  occupied orbitals.
  \item[c)] Derive the factorisation~(\ref{eq:8-factorisation}) from Wick's
  theorem, and verify that the diagonal vanishes at $x_1=x_2$.
  \item[d)] Show that for a correlated state the eigenvalues of
  $\rho^{(1)}$ are strictly between zero and one, and relate this to the
  occupation numbers of Section~\ref{sec:naturalorbitals}.
\end{enumerate}

\paragraph{Answer to a).}
Insert Eq.~(\ref{eq:8-densityoperator}) into the expectation value and use the
fact that $\Psi$ is antisymmetric, so that every term of
$\sum_i\delta(\bm{r}-\bm{r}_i)$ contributes equally:
\begin{equation}
  \bracket{\Psi}{\hat{n}(\bm{r})}{\Psi}
   = N\int\ud{^3\bm{r}_2}\cdots\ud{^3\bm{r}_N}
     \left\vert\Psi(\bm{r},\bm{r}_2,\dots)\right\vert^2 ,
  \label{eq:8-exa1}
\end{equation}
which is Eq.~(\ref{eq:8-onebodydensity}).  Integrating over $\bm{r}$ restores
the full normalisation integral, giving $N\cdot1=N$.

\paragraph{Answer to b).}
For a determinant $\rho^{(1)}(x';x)=\sum_{i=1}^{N}\psi_i(x')\psi_i^{*}(x)$,
which follows from Eq.~(\ref{eq:8-rho1}) by expanding the determinant and
noting that only terms pairing the same orbital in bra and ket survive the
integration over the remaining $N-1$ coordinates.  Setting $x'=x$ gives
Eq.~(\ref{eq:8-densitydeterminant}).  Regarded as an operator,
$\rho^{(1)}=\sum_i\ketbra{\psi_i}{\psi_i}$, so $(\rho^{(1)})^2=\rho^{(1)}$
by orthonormality and $\mathrm{tr}\,\rho^{(1)}=N$: a projector of rank $N$,
which is Exercise~\ref{ex:6-determinants}d) in coordinate language.

\paragraph{Answer to c).}
Write the pair density as
$\bracket{\Phi_0}{\hat{\psi}^{\dagger}(y_1)\hat{\psi}^{\dagger}(y_2)
\hat{\psi}(x_2)\hat{\psi}(x_1)}{\Phi_0}$ and apply Wick's theorem.  Only
contractions pairing a creation with an annihilation operator survive, and
there are exactly two ways to do it: $(y_1x_1)(y_2x_2)$ and
$(y_1x_2)(y_2x_1)$, the second differing from the first by one interchange
and hence carrying a minus sign.  Each contraction is a factor
$\rho^{(1)}$, giving Eq.~(\ref{eq:8-factorisation}).  On the diagonal with
$x_1=x_2$ the two terms become $n(x_1)^2$ and $\vert\rho^{(1)}(x_1;x_1)
\vert^2=n(x_1)^2$, which cancel.

\paragraph{Answer to d).}
$\rho^{(1)}$ is Hermitian and positive semi-definite, so its eigenvalues
$n_k$ are real and non-negative; they sum to $N$.  For fermions no eigenvalue
can exceed one, since $n_k=\bracket{\Psi}{\hat{a}^{\dagger}_k\hat{a}_k}{\Psi}
\le1$ by the Pauli principle.  Equality for exactly $N$ of them forces
$\rho^{(1)}$ to be idempotent, hence forces $\Psi$ to be a single
determinant.  Any correlated state therefore has all $n_k$ strictly inside
$[0,1]$, and the departure of the largest from unity is the diagnostic used
in Table~\ref{tab:natural}.

\subsection*{Exercise 2: the Hohenberg-Kohn theorems}
\label{ex:8-hk}
\addcontentsline{toc}{subsection}{Exercise 2: the Hohenberg-Kohn theorems}

\begin{enumerate}
  \item[a)] Reproduce the proof of the first theorem, stating exactly where
  non-degeneracy is used.
  \item[b)] Show that the theorem determines $v_{\mathrm{ext}}$ only up to an
  additive constant, and explain why that is harmless.
  \item[c)] Prove the second theorem from the first, and state what
  ``universal'' means.
  \item[d)] Explain the $v$-representability problem and how the Levy-Lieb
  constrained search, Eq.~(\ref{eq:8-levylieb}), removes it.
  \item[e)] The theorems say the density determines everything.  Why, then,
  can we not obtain excitation energies from $E[n]$?
\end{enumerate}

\paragraph{Answer to a).}
See Eqs.~(\ref{eq:8-hkineq1})--(\ref{eq:8-hkcontradiction}).  Non-degeneracy
enters in making the variational inequality \emph{strict}: if the ground state
were degenerate, $\Psi^{(2)}$ might be another ground state of
$\hat{H}^{(1)}$ and the inequality would become an equality, so that adding
the two relations would give $E^{(1)}+E^{(2)}\le E^{(1)}+E^{(2)}$ -- true, and
no contradiction.  This is exactly the loophole that the Levy-Lieb
formulation closes.

\paragraph{Answer to b).}
Adding a constant $c$ to $v_{\mathrm{ext}}$ adds $Nc$ to every eigenvalue
and changes no eigenstate, hence no density.  The map from densities to
potentials therefore cannot possibly resolve the constant, and does not need
to: energy differences, densities and all observables are unaffected.  It is
the same ambiguity as the zero of a potential in elementary mechanics.

\paragraph{Answer to c).}
By the first theorem $n_0$ determines $v_{\mathrm{ext}}$, hence $\hat{H}$,
hence its ground state $\Psi[n]$.  So
$F_{\mathrm{HK}}[n]=\bracket{\Psi[n]}{\hat{T}+\hat{V}}{\Psi[n]}$ is a
functional of $n$ alone.  ``Universal'' means that $F$ contains no reference
to $v_{\mathrm{ext}}$: the very same functional applies to a hydrogen atom, a
copper crystal and a stretch of DNA, since $\hat{T}$ and $\hat{V}$ are
properties of electrons.  For any trial $n$, the Rayleigh-Ritz principle
gives $E[n]=F[n]+\int v_{\mathrm{ext}}n \ge E_0$, with equality only at
$n=n_0$.

\paragraph{Answer to d).}
The construction of $F_{\mathrm{HK}}$ requires that $n$ be the ground-state
density of \emph{some} potential in the admissible family -- that it be
$v$-representable.  Densities are known that are not, and no useful
characterisation of the $v$-representable ones exists.  Levy and Lieb
sidestep the issue by defining $F$ through a minimisation over all
antisymmetric wave functions yielding the given density,
Eq.~(\ref{eq:8-levylieb}).  This requires only $N$-representability, which is
easy to characterise for a one-body density -- positivity, correct
normalisation and finite $\int\vert\nabla n^{1/2}\vert^2$ suffice -- and it
handles degenerate ground states.

\paragraph{Answer to e).}
The theorems concern the ground state.  The variational principle
$E[n]\ge E_0$ singles out the minimum and says nothing about stationary
points corresponding to excited states, and the functional $F[n]$ was
constructed from ground-state wave functions only.  Excitations require
extra structure: time-dependent density functional theory obtains them from
the linear response of the density, which is formally the density-functional
analogue of the random-phase approximation of
Chapter~\ref{chap:mfapplications}, with the exchange-correlation kernel
$\delta v_{\mathrm{xc}}/\delta n$ playing the part of the residual
interaction in $A$ and $B$.

\subsection*{Exercise 3: from the electron gas to a functional}
\label{ex:8-gas}
\addcontentsline{toc}{subsection}{Exercise 3: from the electron gas to a functional}

\begin{enumerate}
  \item[a)] Derive the Thomas-Fermi kinetic density~(\ref{eq:8-tf}) from the
  kinetic energy of Chapter~\ref{chap:hartrefock}, and confirm it numerically.
  \item[b)] Derive the local exchange energy~(\ref{eq:8-ldaexchange}) and show
  that it reproduces the $-0.916/r_s$ of Section~\ref{sec:egenergy}.
  \item[c)] Show that $v_x=\tfrac43\epsilon_x$ for the local exchange
  functional, and explain why the factor is not one.
  \item[d)] Derive the one-dimensional analogue~(\ref{eq:8-lda1d}) and the
  kinetic density $t[n]=(\pi^2/6)n^3$.
  \item[e)] Verify the exchange-hole sum rule~(\ref{eq:8-sumrule})
  analytically.
\end{enumerate}

\paragraph{Answer to a).}
Equation~(\ref{eq:6-egkineticdensity}) gives the kinetic energy per particle
as $\tfrac35\hbar^2k_F^2/2m$; multiplying by $n$ and using
$k_F=(3\pi^2n)^{1/3}$ gives the energy per unit volume
\begin{equation}
  t_0[n] = \frac{3}{10}\frac{\hbar^2}{m}\left(3\pi^2\right)^{2/3}n^{5/3},
  \label{eq:8-exa3}
\end{equation}
which is Eq.~(\ref{eq:8-tf}).  The program compares the two expressions at
three densities and finds them identical to eight digits.

\paragraph{Answer to b).}
The exchange energy is
$E_x=-\tfrac12\iint v(\vert\bm{r}-\bm{r}'\vert)
\vert\rho^{(1)}(\bm{r};\bm{r}')\vert^2$, the second term of
Eq.~(\ref{eq:8-pairdiagonal}) integrated against the interaction.  For the
uniform gas $\rho^{(1)}$ depends only on $s=\vert\bm{r}-\bm{r}'\vert$ and
equals $3n[\sin y-y\cos y]/y^3$ with $y=k_Fs$; doing the integral with
$v=e^2/s$ gives $E_x/V=-\tfrac34(3/\pi)^{1/3}n^{4/3}$, hence
Eq.~(\ref{eq:8-ldaexchange}).  Evaluating at $n=3/(4\pi r_0^3)$,
\begin{equation}
  \epsilon_x = -\frac{3}{4}\left(\frac{3}{\pi}\right)^{1/3}
    \left(\frac{3}{4\pi}\right)^{1/3}\frac{1}{r_s}
   = -\frac{0.4582}{r_s}\ \text{Ha}
   = -\frac{0.9163}{r_s}\ \text{Ry},
  \label{eq:8-exb3}
\end{equation}
which is the exchange term of Eq.~(\ref{eq:6-egenergy}).

\paragraph{Answer to c).}
$E_x=\int n\epsilon_x(n)$ with $\epsilon_x\propto n^{1/3}$, so the integrand
is $\propto n^{4/3}$ and
\begin{equation}
  v_x = \frac{\ud{}}{\ud{n}}\left[n\epsilon_x(n)\right]
      = \frac{4}{3}\epsilon_x(n).
  \label{eq:8-exc3}
\end{equation}
The factor is not one because $v_x$ is a functional derivative, not an energy
per particle: changing the density at a point changes both the number of
electrons there and the energy each of them carries.  The distinction is the
same as that between $\sum_i\varepsilon_i$ and the total energy in
Eq.~(\ref{eq:8-kstotalenergy}), and forgetting it is a standard first error.

\paragraph{Answer to d).}
For a spinless one-dimensional gas, $n=\int_{-k_F}^{k_F}\ud{k}/2\pi=k_F/\pi$,
so $k_F=\pi n$, and
$\rho^{(1)}(s)=\int_{-k_F}^{k_F}(\ud{k}/2\pi)e^{iks}=\sin(k_Fs)/\pi s$.
Inserting into the exchange integral and using the symmetry of the two
half-lines to cancel the factor $\tfrac12$ gives Eq.~(\ref{eq:8-lda1d}).  For
the kinetic energy,
$T/L=\int_{-k_F}^{k_F}(\ud{k}/2\pi)(k^2/2)=k_F^3/6\pi=(\pi^2/6)n^3$.

\paragraph{Answer to e).}
Substituting Eq.~(\ref{eq:8-paircorrelation}),
\begin{equation}
  \frac{4}{3\pi}\int_0^{\infty}\ud{y}\,y^2\left[g(y)-1\right]
   = -\frac{4}{3\pi}\cdot\frac{9}{2}
     \int_0^{\infty}\ud{y}\,\frac{(\sin y-y\cos y)^2}{y^4}
   = -\frac{6}{\pi}\cdot\frac{\pi}{6} = -1,
  \label{eq:8-exe3}
\end{equation}
using the standard integral
$\int_0^{\infty}(\sin y-y\cos y)^2y^{-4}\ud{y}=\pi/6$.  The program obtains
$-0.9976$ by direct quadrature, the residue being the slowly decaying
oscillatory tail.

\subsection*{Exercise 4: the Kohn-Sham equations}
\label{ex:8-kohnsham}
\addcontentsline{toc}{subsection}{Exercise 4: the Kohn-Sham equations}

\begin{enumerate}
  \item[a)] Derive Eq.~(\ref{eq:8-kohnsham}) by varying
  Eq.~(\ref{eq:8-ksenergy}) with respect to the orbitals subject to
  orthonormality.
  \item[b)] Derive the total-energy expression~(\ref{eq:8-kstotalenergy}) and
  explain each correction term.
  \item[c)] Compare Eqs.~(\ref{eq:8-hfcoordinate})
  and~(\ref{eq:8-kscoordinate}) term by term.  Which is variational, and with
  respect to what?
  \item[d)] Why is $T_s[n]$, and not $T[n]$, computed from the orbitals?
  What is the difference, and where does it go?
  \item[e)] Show that for one electron the exact $E_{\mathrm{xc}}$ must
  satisfy $E_{\mathrm{xc}}[n]=-\mathcal{U}_{\mathrm{H}}[n]$.
\end{enumerate}

\paragraph{Answer to a).}
Write $T_s$ and $n$ in terms of the orbitals and impose
$\int\psi_i^{*}\psi_j=\delta_{ij}$ with multipliers $\varepsilon_{ij}$.
Varying with respect to $\psi_i^{*}$,
\begin{equation}
  -\frac{\hbar^2}{2m}\nabla^2\psi_i
  + \left[v_{\mathrm{ext}}
    + \frac{\delta\mathcal{U}_{\mathrm{H}}}{\delta n}
    + \frac{\delta E_{\mathrm{xc}}}{\delta n}\right]
    \frac{\delta n}{\delta\vert\psi_i\vert^2}\psi_i
  = \sum_j\varepsilon_{ij}\psi_j ,
  \label{eq:8-exa4}
\end{equation}
and since $\delta n/\delta\vert\psi_i\vert^2=1$ the bracket is exactly
$v_{\mathrm{KS}}$.  A unitary rotation among the occupied orbitals
diagonalises $\varepsilon_{ij}$ without changing the density -- the same
freedom as in Section~\ref{sec:detbasis} -- giving
Eq.~(\ref{eq:8-kohnsham}).

\paragraph{Answer to b).}
Multiplying Eq.~(\ref{eq:8-kohnsham}) by $\psi_i^{*}$, integrating and
summing over the occupied orbitals gives
$\sum_i\varepsilon_i=T_s+\int v_{\mathrm{KS}}n$.  Substituting
$v_{\mathrm{KS}}$ from Eq.~(\ref{eq:8-kspotential}) and comparing with
Eq.~(\ref{eq:8-ksenergy}),
\begin{equation}
  E = \sum_i\varepsilon_i
    - \underbrace{\int v_{\mathrm{H}}n}_{=2\mathcal{U}_{\mathrm{H}}}
    - \int v_{\mathrm{xc}}n
    + \mathcal{U}_{\mathrm{H}} + E_{\mathrm{xc}} ,
  \label{eq:8-exb4}
\end{equation}
which is Eq.~(\ref{eq:8-kstotalenergy}).  The band sum counts the Hartree
energy twice, once for each partner, so one copy is removed; and it contains
$\int v_{\mathrm{xc}}n$, which is not $E_{\mathrm{xc}}$, so the former is
subtracted and the latter added.  This is the same double-counting
correction as Exercise~\ref{ex:6-hfequations}b).

\paragraph{Answer to c).}
The kinetic, external and Hartree terms are identical.  The difference is the
last term: Hartree-Fock has a non-local exchange \emph{operator}, acting
differently on each orbital, exact for a determinant and containing no
correlation; Kohn-Sham has a local multiplicative \emph{potential}, the same
for all orbitals, approximate, but containing correlation as well.
Hartree-Fock is variational with respect to the true energy -- it is a
genuine upper bound on $E_0$.  Kohn-Sham is variational only with respect to
the energy of the model functional it is given; with an approximate
$E_{\mathrm{xc}}$ the result may lie above or below the exact energy, as
Table~\ref{tab:ksresults} shows.

\paragraph{Answer to d).}
$T[n]$, the true kinetic energy of the interacting system, is not accessible
from the density; that was the problem in the first place.  $T_s[n]$, the
kinetic energy of the non-interacting system with the same density, is
computed exactly from its orbitals.  The difference $T[n]-T_s[n]$ is small --
it is a correlation effect -- and it is folded into $E_{\mathrm{xc}}$ by the
definition~(\ref{eq:8-excdefinition}).  The whole construction is the
observation that this is a far better place to put the difficulty than the
kinetic energy itself.

\paragraph{Answer to e).}
For $N=1$ the exact energy contains no electron-electron interaction at all.
But Eq.~(\ref{eq:8-ksenergy}) contains $\mathcal{U}_{\mathrm{H}}[n]$, which is
non-zero for any density.  The exact functional must therefore cancel it
exactly, $E_{\mathrm{xc}}[n]=-\mathcal{U}_{\mathrm{H}}[n]$ for any one-electron
density.  The local density approximation violates this: the program finds a
residue $0.0469$ for the harmonic trap, Eq.~(\ref{eq:8-sie}).  Imposing the
one-electron condition exactly is the content of the Perdew-Zunger
self-interaction correction.

\subsection*{Exercise 5: numerical explorations}
\label{ex:8-numerical}
\addcontentsline{toc}{subsection}{Exercise 5: numerical explorations}

\begin{enumerate}
  \item[a)] Using \texttt{dft.py}, reproduce Table~\ref{tab:ksresults} and
  plot the Kohn-Sham, Hartree-Fock and non-interacting densities for $N=4$.
  Where do they differ most?
  \item[b)] Vary the interaction strength and follow the ratio
  $E_x^{\mathrm{local}}/E_x^{\mathrm{exact}}$.  Does the local approximation
  improve or deteriorate as the interaction is made stronger?
  \item[c)] Implement the Perdew-Zunger self-interaction correction for
  $N=1$ -- subtract $\mathcal{U}_{\mathrm{H}}[n_i]+E_x[n_i]$ orbital by
  orbital -- and check that the one-particle energy returns to $0.5$.
  \item[d)] Solve the Thomas-Fermi equation for the trap directly by
  minimising over $n(x)$ on a grid, and compare the resulting density with
  the Kohn-Sham one.  Is the shell structure visible?
  \item[e)] Add a crude gradient correction,
  $E_x\to E_x-\beta\int\vert\nabla n\vert^2/n^{4/3}$, and tune $\beta$ to fit
  the exact exchange energies of Section~\ref{sec:howlocal}.  Does one value
  work for all $N$?
\end{enumerate}

\paragraph{Answer.}
Some pointers.  In a) the densities differ most in the tails, where the
Kohn-Sham density is too diffuse -- the signature of self-interaction, since a
spuriously self-repelled electron spreads out.  In b) the ratio deteriorates
as the interaction strengthens, because exchange grows relative to the
kinetic energy and the errors in it become more visible; this is the same
trend that made the mean field fail at strong coupling in
Section~\ref{sec:fcihubbard}.  In c) the correction is exact by construction
for one particle and the energy returns to $0.5$ to machine precision; the
interesting question is what it does for $N>1$, where it is no longer exact
and can over-correct.  In d) the Thomas-Fermi density is a smooth dome with
no oscillations, in contrast to the $N$ visible maxima of the Kohn-Sham
density -- the clearest single illustration of what orbitals buy.  In e) a
single $\beta$ does not fit all $N$; the failure of a one-parameter gradient
correction is precisely why the construction of gradient functionals is done
by imposing exact constraints instead of fitting.
